The checked source pin is \path{29936089835c0c55cca81edc86b4fb0c70ce34d5},
using Lean 4.32.0 (compiler commit \path{8c9756b28d64dab099da31a4c09229a9e6a2ef35})
and mathlib \path{81a5d257c8e410db227a6665ed08f64fea08e997}.
The relevant project dependency closure was rebuilt from source: 62 project modules
and one otherwise missing mathlib module. The remaining mathlib and package objects
were reused from a cache matching the pinned dependency manifest. All 63 recorded
compilations completed successfully and produced nonempty objects.

The archive's own replay script then completed the following checks:
\begin{itemize}
\item All four complete proposed production modules compiled successfully in a fresh
object overlay of that baseline.
\item The five changed theorem headers matched their original headers exactly, and
baseline and replacement \code{\#print axioms} reports were identical.
\item The existing \code{WangTiles942HelicalQuotientChecks} and
\code{RatioMarginFinalizer} clients compiled against the proposed modules.
\item The two optional modules, \code{SignedPrefix} and \code{RankPotentialActual},
compiled against the actual production definitions. The signed-prefix file checks
a negative total and an empty prefix; the rank file includes the equal-rank family
and the degenerate base-one case.
\item The separate helical and arithmetic proof examples also compiled. No checked
axiom report contains \code{sorryAx}. The extra helical endpoint audit retained its
original axiom lists as well.
\end{itemize}

The periodic-orbit and coprime-divisibility results depend on \code{propext} and
\code{Quot.sound}. The geometric and margin results also use
\code{Classical.choice}; these are exactly their original axiom lists.
The full command receipts, source hashes, output logs, and line counts are included.
Recorded elapsed times describe these runs only and are not before/after performance
measurements. Validation is against the stated source pin, not against the newer
freshness-check commit described in Section~1.
